\wordchapter{2}{Institutional Background and Literature Review}

\wordsection{2.1}{Program Structure and Premium Components}

MAFRA delegates management of the program to the Agricultural Policy Insurance \& Finance Service (APFS). NH Nonghyup Property \& Casualty Insurance (NH) is the sole operator. It sells contracts through local agricultural cooperatives, services policies, and pays indemnities. Figure 1 shows the payments and risk transfers among these parties.

\begin{wordfigure}{1}{Organizational Structure of Crop Disaster Insurance in South Korea}
\IfFileExists{"organization.tex"}{%
  {\centering\adjustbox{max width=6.2in}{\input{"organization"}}\par}%
}{%
  \wordimage{6.2in}{figures/fig1_organization.png}%
}
\wordnotes{\emph{Notes}: The arrows show payments, and the report of the loss assessment results; the dashed lines represent verification and approval of premium rates and policy terms. Local agricultural cooperatives are not displayed, although they offer contracts for sale on behalf of the operator. Adapted from MAFRA (2025, 2026a).}
\end{wordfigure}

NH transfers part of its risk to private reinsurers and the Agricultural and Fishery Disaster Reinsurance Fund. In 2024, the first-stage shares under the main program's profit-and-loss-sharing arrangement were 10 percent for NH, 40 percent for private reinsurers, and 50 percent for the state (MAFRA, 2025, p. 274). Additional non-proportional provisions alter the allocation at higher loss levels. These initial shares therefore do not describe the allocation of every layer of loss.

The gross premium consists of a pure premium and a loading. The pure premium covers insured risk and loss-investigation expenses. The loading covers specified operating expenses, including delegated selling costs. Under the 2026 rules, the government finances part of the pure premium and the entire loading (MAFRA, 2026a). The subsidy rate in equation (1) applies only to the pure premium. Public support consequently accounts for a larger share of the gross premium than this rate alone indicates.

\wordsection{2.2}{Expansion, Enrollment, and Financing}

The program began with apples and pears in 2001 and expanded to 78 crops and 94 products by 2026 (MAFRA, 2026a). A crop's pilot, expansion of sales regions, and administrative entry into the main program need not coincide. In the national analysis, I use the last pre-pilot year as the reference and examine ten target years beginning with the first harvest covered by documented nationwide sales (Table 3). These dates measure availability, while enrollment measures farmers' use of that availability.

In the series used for Figure 2, the national enrollment rate rose from 12.5 percent in 2009 to 54.2 percent in 2024 and stood at 54.0 percent in 2025, when insured area declined from 676,226 to 655,456 hectares (APFS, n.d.-a; MAFRA, 2026b). This rate divides insured area by eligible area. Changes in the rate reflect both take-up and changes in eligibility as crops or regions enter the program. The loss ratio divides portfolio claims by the premium measure specified in the figure notes. It measures claims relative to premiums, rather than individual farms' income losses or NH's company-wide profitability after reinsurance.

Figure 3 reports the financing shares of the pure premium in the APFS series used here. In 2024, local governments paid 39.7 percent, policyholders paid 12.6 percent, and the central government paid 47.7 percent (APFS, n.d.-a). These shares exclude the publicly financed loading. Transfers to individual policyholders depend on insured value, rated risk, and coverage.

\begin{wordfigure}{2}{National Enrollment Rate and Loss Ratio, 2001--2025}
\wordimage{5.27746in}{figures/fig2_enrollment_loss_ratio.png}
\wordnotes{\emph{Notes:} The enrollment is the ratio of insured area to eligible area. Loss ratio is the ratio of indemnities to premium net of refunds for 2011 and risk premium for 2012. The 100 percent horizontal line is shown. Enrollment through 2024 is from APFS (n.d.-a), and loss ratios through 2023 from APFS (n.d.-b). The 2025 enrollment rate and 2024--2025 loss ratios are from MAFRA (2026b, p. 4).}
\end{wordfigure}

Enrollment varies substantially across crops (Table 1). In 2023, insurance covered 89.0 percent of eligible apple area and 61.0 percent of eligible rice area, compared with 7.1 percent for sweet potato and 5.4 percent for corn. Insurance availability thus corresponded to very different levels of coverage among the six annual crops examined in Chapter 5.

\begin{wordfigure}{3}{Shares of the Pure Premium by Payer, 2001-2024}
\wordimage{5.4in}{figures/fig3_pure_premium_shares.png}
\wordnotes{\emph{Notes}: Shares of the reported pure-premium total. Unattributed segments in 2010-2011 are the difference between that total and the identified payer components. The publicly financed loading is excluded. Data from APFS (n.d.-a).}
\end{wordfigure}

\begin{wordtable}{1}{Enrollment and Loss Ratios for Selected Crops, 2023}
\begin{wordtabular}{L{0.304}C{0.185}C{0.182}C{0.173}C{0.156}}
\wtop
\hdl{0.304}{Crop} & \hd{0.185}{Eligible area (ha)} & \hd{0.182}{Insured area (ha)} & \hd{0.173}{Enrollment (\%)} & \hd{0.156}{Loss ratio (\%)} \\
\whead
\panelgap
\textit{Panel A. Rice and major fruits} & & & & \\
Rice & 725,882 & 443,136 & 61.0 & 103.5 \\
Apple & 24,945 & 22,209 & 89.0 & 116.1 \\
Pear & 10,078 & 7,683 & 76.2 & 60.5 \\
\panelgap
\textit{Panel B. Six annual crops} & & & & \\
Soybean & 55,326 & 28,649 & 51.8 & 158.9 \\
Onion & 17,128 & 6,811 & 39.8 & 73.1 \\
Garlic & 21,110 & 5,936 & 28.1 & 119.0 \\
Red pepper & 26,126 & 8,522 & 32.6 & 141.7 \\
Sweet potato & 18,389 & 1,300 & 7.1 & 195.5 \\
Corn & 14,073 & 764 & 5.4 & 129.2 \\
\wbottom
\end{wordtabular}

\wordnotes{\emph{Notes}: The area eligible may vary from the national cultivated area due to the sales areas of each product. Loss ratio is the ratio of indemnities to risk premium. Data from APFS (n.d.-c); calculations made by the author.}
\end{wordtable}

\wordsection{2.3}{Related Literature}

The literature on agriculture public support is motivated by two features of agricultural insurance. First, insurers observe farm risks and farmers' actions imperfectly. Nelson and Loehman (1987) examine the implications of adverse selection and moral hazard for contract design. Just et al. (1999) show that participation reflects both subsidy incentives and asymmetric information, while Lee and Jung (2014) examine these information problems among Korean rice farms. Second, weather shocks can generate losses across many policyholders at once. Miranda and Glauber (1997) examine the resulting systemic risk and the role of reinsurance.

The next question is how farmers, insurers, and the government should share costs and risk. Jia et al. (2025) show theoretically that the public role depends on loss correlation and the relative efficiency of public and private risk bearing. Kang and Chung (2017) evaluate Korea's risk-group reinsurance arrangement piloted during 2014--2016 and its implications for losses and default risk in the state fund. For premium subsidies, Hazell et al. (2017) emphasize explicit objectives, targeting, and cost-effective design. Insurance demand may also depend on expected assistance outside the insurance market (Raschky \& Weck-Hannemann, 2007). Participation alone therefore provides an incomplete measure of performance. Lusk (2017) examines the distribution of subsidy benefits and costs among producers, consumers, and taxpayers.

Insurance may also have an impact on production through altering the expected income and risk. Goodwin et al. (2004) find modest acreage responses in the U.S. crop regions they examine, and Yu et al. (2018) estimate acreage responses to premium subsidies. In Korea, Han (2014) compares insurable and non-insurable crops to examine production and price effects. Nam and Ahn (2022) study input use and farm income, while Lee and Kim (2026) examine output and technical efficiency using fruit-farm panel data. These studies provide evidence on production decisions and farm outcomes. The present study places that evidence alongside an assessment of the purposes and incentives of public support.

There are also methodological issues regarding timing of insurance introduction. Callaway and Sant'Anna (2021) develop a framework for comparisons across treatment cohorts, including comparisons with units not yet treated. Roth (2022) shows why a statistically insignificant pre-trend test should not be treated as proof of parallel trends. The literature on clustered inference addresses uncertainty when observations are dependent and the number of treated clusters is small (Cameron et al., 2008; MacKinnon \& Webb, 2018; MacKinnon et al., 2023; Webb, 2023). These studies inform the choice of comparison groups and uncertainty measures used here.

I bring these perspectives together by assessing each support instrument and comparing changes in area, yield, and production around insurance introduction. The national analysis is done on a crop-by-crop basis. Crop identity is kept as constant in the rice analysis to explore the first two years of availability in municipalities. The two comparisons convey production responses to the package, and the institutional assessment signifies the other evidence desired to assess the parts of the policy package.
